\begin{exercise}[subtitle={Drunk Men and Drunk Birds}]

There is a saying, attributed to  Japanese-American mathematician Shizuo
Kakutani: ``A drunk man will find his way home, but a drunk bird may get lost
forever." Let us try to understand this. A drunk birds can be modeled by three
dimensional random walk. Suppose the bird starts at the origin of the three
dimensional integer lattice $\integer^3$. It moves only at integer time points. For
$n=0,1,2,3,\ldots$ let $S_n=(S_n^x,S_n^y,S_n^z)\in\integer^3$ denote its location.
If at a certain time the bird is at location $(x,y,z)$, then at the next time
point it can be at one of the $8$ locations given by $(x\pm 1,y\pm 1,z\pm 1)$.
Thus we can say

\begin{align*}
S_n^x ={} & S_{n-1}^x + X_n\;,\\
S_n^y ={} & S_{n-1}^y + Y_n\;,\\
S_n^z ={} & S_{n-1}^z + Z_n\;,
\end{align*}
where
\[
X_n = \begin{cases}
1 & \mbox{ with probability } 1/2\;,\\
-1 & \mbox{ with probability } 1/2\;,
\end{cases}
\]
\[
Y_n = \begin{cases}
1 & \mbox{ with probability } 1/2\;,\\
-1 & \mbox{ with probability } 1/2\;,
\end{cases}
\]
\[
Z_n = \begin{cases}
1 & \mbox{ with probability } 1/2\;,\\
-1 & \mbox{ with probability } 1/2\;,
\end{cases}
\]
and $(X_1,X_2,X_3,\ldots,Y_1,Y_2,Y_3,\ldots,Z_1,Z_2,Z_3,\ldots)$ are independent.
\[
\prob*{S_n=(0,0,0)} 
= \prob*{S_n^x=0}\prob*{S_n^y=0}\prob*{S_n^z=0}
\]
For $n$ odd, 
\[
\prob*{S_n^x=0}=
\prob*{S_n^y=0}=
\prob*{S_n^z=0}=0\;.
\]
For $n$ even, say $n=2N$
\[
\prob*{S_n^x=0}=
\prob*{S_n^y=0}=
\prob*{S_n^z=0}=\left(2^{-2N}\binom{2N}{N}\right)^3\;.
\]
By Stirling's approximation
\[
2^{-2N}\binom{2N}{N} \approx \frac{C}{N^{1/2}} 
\approx \frac{C}{n^{1/2}}.
\]
Thus, for $n$ even
\[
\prob*{S_n=0} \approx C n^{-3/2}\;.
\]
Thus, the probability that $S_n=0$ happens infinitely often is $0$. Thus, a
drunk bird may get lost forever. But for two dimensional or one dimensional
random walk with probability $1$ $S_n=0$ happens infinitely often. Thus a drunk
man will find his way home. Note that, to prove this we cannot use the second
Borel-Cantelli Lemma. Because the second Borel-Cantelli Lemma requires the
events to be independent, but the events $\{S_n=0\}$ are not independent. But
the result is true and can be proved in different ways, generally you will learn
this in a course on Markov Chains.

\end{exercise}
